% TOP_PROBLEM: {"id": 2864, "title": "Kirby Problem 3.66", "category_id": 7, "category": {"id": 7, "name": "topology", "display_name": "Topology", "description": "Properties preserved under continuous deformations.", "slug": "topology", "order_index": 7, "created_at": "2026-07-31T15:26:25.671Z"}, "status": "open"}
% FIRST_POSTED: null
% ENTRY_KIND: statement_only
% Imported from: batch12.tex; UnsolvedMath ID 2864
\documentclass[11pt]{article}
\usepackage[margin=25mm]{geometry}
\usepackage{fontspec}
\setmainfont{Latin Modern Roman}
\usepackage{amsmath,amssymb,mathtools}
\usepackage{unicode-math}
\setmathfont{Latin Modern Math}
\usepackage{xurl,hyperref}
\hypersetup{colorlinks=true,urlcolor=blue}
\setcounter{secnumdepth}{0}
\setlength{\parindent}{0pt}
\setlength{\parskip}{5pt}
\setlength{\emergencystretch}{3em}
\newcommand{\sref}[2]{\hyperlink{source:#1:#2}{[S#2]}}
\newcommand{\cataloguescope}[1]{\paragraph{Recorded target.}#1}
\begin{document}
% UnsolvedMath ID 2864; source catalogue code KP-3.66
\setcounter{section}{2863}
\section{Kirby Problem 3.66}\label{2864}
{\small\sffamily UnsolvedMath ID 2864 · Topology}
\subsection{Definitions and mathematical statement}

\paragraph{Shared notation.}$\mathbb N=\{1,2,\ldots\}$, and $\mathbb Z,\mathbb Q,\mathbb R,\mathbb C$ denote the integers, rationals, reals and complex numbers. Any other conventions are specified locally.


Put $R=\mathbb Z[A,A^{-1}]$ and $R_{\mathbb Q}=\mathbb Q[A,A^{-1}]$. For a closed oriented $3$-manifold $Y$, the Kauffman bracket skein module $S(Y)$ is the free $R$-module on isotopy classes of unoriented framed links, including the empty link, modulo the local bracket relations
\[
[L_\times]=A[L_0]+A^{-1}[L_\infty],\qquad
[L\sqcup O]=-(A^2+A^{-2})[L].
\]
The crossing and its two smoothings are the usual ordered Kauffman bracket local pictures in a ball, and $O$ is a disjoint zero-framed trivial circle. Let $S_{\mathbb Q}(Y)=S(Y)\otimes_R R_{\mathbb Q}$. The source calls this module tame if it is a direct sum of cyclic $R_{\mathbb Q}$-modules and there exists an odd integer $N\ge1$ for which it has no submodule isomorphic to $R_{\mathbb Q}/(\Phi_{2N}(A))$, where $\Phi_{2N}$ is the cyclotomic polynomial.
Let $\operatorname{HP}^*_{\#}(Y)$ be Abouzaid--Manolescu's framed sheaf-theoretic $\operatorname{SL}_2(\mathbb C)$ Floer cohomology: it is the hypercohomology of the perverse sheaf of vanishing cycles for the framed complex-Lagrangian intersection associated to a Heegaard splitting, on $\operatorname{Hom}(\pi_1(Y),\operatorname{SL}_2(\mathbb C))$. Its degree-zero group is $\operatorname{HP}^0_{\#}(Y)$; take its dimension after tensoring its integral coefficients with $\mathbb C$. \sref{2864}{3}
\paragraph{Statement recorded in the supplied source.}
For every closed oriented $3$-manifold $Y$:
\begin{enumerate}
\item Is
\[
\dim_{\mathbb Q(A)}\bigl(S(Y)\otimes_R\mathbb Q(A)\bigr)
=\dim_{\mathbb C}\bigl(\operatorname{HP}^0_{\#}(Y)\otimes_{\mathbb Z}\mathbb C\bigr)?
\]
\item Characterize exactly when $S_{\mathbb Q}(Y)$ is tame in the sense above.
\end{enumerate}
\sref{2864}{2}
\subsection{Short English statement}
Does generic Kauffman bracket skein rank equal degree-zero framed complex Floer dimension for every closed oriented three-manifold, and which Laurent-polynomial skein modules are tame?
\subsection{Sources}
\begin{description}
\item[\hypertarget{source:2864:1}{S1}] ULAM.AI and the UnsolvedMath Contributors, \emph{UnsolvedMath: A Curated Collection of Open Mathematics Problems}, record 2864 (KP-3.66). CC BY 4.0. \url{https://huggingface.co/datasets/ulamai/UnsolvedMath}.
\item[\hypertarget{source:2864:2}{S2}] R. İnanç Baykur, Robion C. Kirby and Daniel Ruberman (editors), \emph{K3: A New Problem List in Low-Dimensional Topology}, Mathematical Surveys and Monographs 295 (2026), Problem 3.66, p. 179 and following remarks; original author version. \url{https://math.berkeley.edu/sites/default/files/surv-295-ruberman-watermarked-author-pdf.pdf}.
\item[\hypertarget{source:2864:3}{S3}] Mohammed Abouzaid and Ciprian Manolescu, \emph{A sheaf-theoretic model for SL(2,C) Floer homology}, Journal of the European Mathematical Society 22 (2020), 3641--3695; arXiv:1708.00289v2, introduction, Theorem 1.3 and the framed hypercohomology definition. \url{https://arxiv.org/abs/1708.00289v2}.
\end{description}
\label{end:2864}
\end{document}
